\section{Certificate, Route-A status and limitations}

The finite certificate is deliberately smaller than the source theorem.  It
locks eight inherited artifacts, reconstructs all reciprocal trace roots for
the three exact orbits, verifies the residue interval and the primitive versus
repetition logarithmic-derivative identity, and rejects twelve adversarial
mutations in an independent implementation.  Ten unit tests cover the same
interfaces.  The complete finite audit is reproduced by
\texttt{bash code/run\_c54.sh} from the project directory.

The exact rows have two roles.  They ensure that no code path silently equates
Mahler height with the physical multiplier, and they prove the scalar-roof
obstruction.  They do not establish a H\"older potential or estimate
\(\sigma_{\rm Gal}\).

For Route A, the normalized physical suspension already has an analytic
dynamical zeta, while the full Galois-weighted amplitude has no known Euler or
Fredholm determinant.  We therefore report the coordinatewise assessment
\[
\begin{aligned}
(&\texttt{A1\_WEAK},
\ \texttt{A2\_ANALYTIC\_DETERMINANT}
   \text{ (physical subsystem)},\\
 &\texttt{A3\_PARTIAL\_ANALYTIC\_STRUCTURE},
\ \texttt{A4\_FORMAL\_HINT}),
\end{aligned}
\]
with overall status \texttt{ROUTE\_A\_EXPLORATORY}.  The word ``physical'' is
part of the A2 scope.  For the full Galois-weighted candidate, the dynamical
determinant interface remains open and no A2 pass is claimed.
Route B is not authorized.

The strongest positive result is the exact pressure pole of the physical
Mahler summand.  The strongest obstruction is the failure of scalar pressure
retuning and the absence of any exact symbolic realization of the Galois
excess.  The next theorem should either construct a H\"older or controlled
asymptotically additive excess observable, or compute the excess abscissa
\(\sigma_{\rm Gal}\) directly.

Nothing here supplies rational-prime labels, von Mangoldt amplitudes, a
functional equation, a completed Riemann divisor, or a self-adjoint operator.
